\par\bigskip\Needspace{9\baselineskip}
\pdfbookmark[2]{Zweite Methode, [5]}{ceres-second-5}
\begin{center}[5.]\end{center}
\begingroup\small
\noindent\hspace*{\parindent}[Nachdem auf diese Weise alle Glieder des Ausdrucks 2) auf die Form
\[
A\sin(iD\pm\text{\AstroSymbols ⚳}\mp\AscendingNode+kM-k'M')
\]
gebracht sind, wird das entsprechende
\[
\begin{aligned}
&5)\qquad\text{Glied in }\frac{\theta}{\sin i}=\\
&\quad-A\cos V\int\sin V.\sin(iD\pm\text{\AstroSymbols ⚳}\mp\AscendingNode+kM-k'M')\,n\,\mathrm{d}t\\
&\quad+A\sin V\int\cos V.\sin(iD\pm\text{\AstroSymbols ⚳}\mp\AscendingNode+kM-k'M')\,n\,\mathrm{d}t.
\end{aligned}
\]

Der nächste Schritt besteht nun in der Entwickelung von \(\displaystyle\genfrac{}{}{0pt}{}{\sin}{\cos}V\) nach Vielfachen von \(M\), welche auch nach der interpolatorischen Methode erfolgen kann. Es fand sich
\[
\begin{aligned}
\cos V&=0{,}07848+0{,}99308\cos M-0{,}07784\cos2M\quad\text{etc.}\\
\sin V&=\phantom{0{,}07848+}0{,}99461\sin M-0{,}07792\sin2M\quad\text{etc.}
\end{aligned}
\]

Nachdem die Ausmultiplication mit diesen Grössen unter den Integralzeichen in 5) ausgeführt ist, kann die Ausführung der Integration ohne Weiteres erfolgen und nach nochmaliger Multiplication mit \(\cos V\) resp. \(\sin V\) erhält man das entsprechende Glied in \(\theta\). Das oben als Beispiel benutzte Glied 3) ergab so:
\SourceLeaf{413}
\[
\begin{aligned}
&-27''{,}47\sin(3D+\text{\AstroSymbols ⚳}-\AscendingNode)\\
&-\phantom{2}4''{,}56\sin(3D+\text{\AstroSymbols ⚳}-\AscendingNode+M)\\
&+\phantom{2}2''{,}26\sin(3D+\text{\AstroSymbols ⚳}-\AscendingNode-M)\\
&+\phantom{2}2''{,}81\sin(3D+\text{\AstroSymbols ⚳}-\AscendingNode+M').
\end{aligned}
\]

Die gesammten numerischen Resultate finden sich im folgenden Briefe, aus welchem auch zu ersehen ist, wie die Glieder, deren Argumente sich nur um constante Grössen unterscheiden, in eines zusammengezogen worden sind.]
\par\endgroup
\begin{center}\rule{.30\linewidth}{.4pt}\end{center}
